% Catalogo lexical: tabelas de DEFINICAO dos campos figurativos.
%
% ATENCAO: estas sao tabelas de CATALOGO (definicao), nao de frequencia.
% Listam QUAIS TERMOS compoem cada campo, na ordem do arquivo de definicao,
% independentemente de contagem ou obra. Nao confundir com
% outputs/latex/campos_lexicais_detalhados.tex, que traz as FREQUENCIAS por obra.
%
% Versao enxuta: campos com ate 8 termos aparecem completos; campos com mais
% de 8 mostram os 8 primeiros termos do arquivo de definicao seguidos de
% "entre outros". A ordem dos termos e a ordem do arquivo de definicao.
%
% Pacotes necessarios: \usepackage{booktabs, longtable, array}.
%
% Fontes da verdade (catalogo completo no repositorio):
%   Tabela A (Etapas 1-2): campos_lexicais/catalogo_termos.yaml
%       + campos_lexicais/latour_textil_en_etapa2_adicoes.txt
%       + campos_lexicais/latour_topologia_en_etapa2_adicoes.txt
%   Tabela B (Etapa 3): campos_lexicais/catalogo_termos_aime.yaml

\begin{longtable}{@{}p{0.23\textwidth} p{0.20\textwidth} p{0.49\textwidth}@{}}
\caption{Catalogo lexical das Etapas 1 e 2: os 19 campos figurativos rastreados nas obras de Latour, com a procedencia de cada campo e uma amostra dos termos que o definem. Amostra do catalogo: campos com ate oito termos aparecem completos, campos maiores trazem os oito primeiros termos do arquivo de definicao seguidos de \enquote{entre outros}, sempre na ordem do arquivo. A lista completa de termos esta no repositorio, em \texttt{campos\_lexicais/catalogo\_termos.yaml} (17 campos) e nos arquivos de adicao da Etapa 2 \texttt{campos\_lexicais/latour\_textil\_en\_etapa2\_adicoes.txt} e \texttt{campos\_lexicais/latour\_topologia\_en\_etapa2\_adicoes.txt}.}\label{tab:catalogo-etapas12}\\
\toprule
Campo & Procedencia & Amostra de termos do catalogo \\
\midrule
\endfirsthead
\caption[]{Catalogo lexical das Etapas 1 e 2 (continuacao).}\\
\toprule
Campo & Procedencia & Amostra de termos do catalogo \\
\midrule
\endhead
\midrule
\multicolumn{3}{r@{}}{\footnotesize continua na proxima pagina}\\
\endfoot
\bottomrule
\endlastfoot
\texttt{inscription} & conceito de Latour & \texttt{inscription}, \texttt{inscriptions}, \texttt{inscription device}, \texttt{inscription devices}, \texttt{literary inscription} \\
\texttt{immutable\_mobile} & conceito de Latour & \texttt{immutable mobile}, \texttt{immutable mobiles} \\
\texttt{black\_box} & conceito de Latour & \texttt{black box}, \texttt{black-box}, \texttt{black boxes}, \texttt{blackbox}, \texttt{blackboxing}, \texttt{black-boxing} \\
\texttt{centre\_of\_calculation} & conceito de Latour & \texttt{centre of calculation}, \texttt{center of calculation}, \texttt{centres of calculation}, \texttt{centers of calculation} \\
\texttt{actor\_network} & conceito de Latour & \texttt{actor-network}, \texttt{actor network}, \texttt{actant}, \texttt{actants} \\
\texttt{translation} & conceito de Latour & \texttt{translation}, \texttt{translations}, \texttt{translate}, \texttt{translates}, \texttt{translated} \\
\texttt{trial\_of\_strength} & conceito de Latour & \texttt{trial of strength}, \texttt{trials of strength} \\
\texttt{factish} & conceito de Latour & \texttt{factish}, \texttt{factishes} \\
\texttt{circulating\_reference} & conceito de Latour & \texttt{circulating reference}, \texttt{circulating references} \\
\texttt{articulation} & conceito de Latour & \texttt{articulation}, \texttt{articulations}, \texttt{articulate}, \texttt{articulated} \\
\texttt{construction} & conceito de Latour & \texttt{construction}, \texttt{constructed}, \texttt{constructing}, \texttt{social construction} \\
\texttt{proposition} & conceito de Latour & \texttt{proposition}, \texttt{propositions} \\
\texttt{network} & conceito de Latour & \texttt{network}, \texttt{networks}, \texttt{networking} \\
\texttt{agonistic} & conceito de Latour & \texttt{agonistic}, \texttt{agonistics}, \texttt{agonistic field} \\
\texttt{enrollment} & conceito de Latour & \texttt{enrollment}, \texttt{enrolment}, \texttt{enroll}, \texttt{enrol}, \texttt{enrolled} \\
\texttt{spokesperson} & conceito de Latour & \texttt{spokesperson}, \texttt{spokespersons}, \texttt{spokesman}, \texttt{spokesmen}, \texttt{spokeswoman} \\
\texttt{militar} & campo reunido & \texttt{ally}, \texttt{allies}, \texttt{allied}, \texttt{alliance}, \texttt{alliances}, \texttt{enemy}, \texttt{enemies}, \texttt{enmity}, entre outros \\
\texttt{textil} & campo reunido & \texttt{thread}, \texttt{threads}, \texttt{threading}, \texttt{threaded}, \texttt{thread-like}, \texttt{threadlike}, \texttt{weave}, \texttt{weaves}, entre outros \\
\texttt{topologia} & campo reunido & \texttt{fluid}, \texttt{fluids}, \texttt{fluidity}, \texttt{surface}, \texttt{surfaces}, \texttt{node}, \texttt{nodes}, \texttt{connection}, entre outros \\
\end{longtable}

\begin{longtable}{@{}p{0.23\textwidth} p{0.23\textwidth} p{0.46\textwidth}@{}}
\caption{Catalogo lexical da Etapa 3 (AIME, Latour 2013): os 12 campos especificos identificados a partir do Quadro 1 e da Conclusao de \emph{An Inquiry into Modes of Existence}, com a procedencia e uma amostra dos termos que definem cada campo. Amostra do catalogo: campos com ate oito termos aparecem completos, campos maiores trazem os oito primeiros termos do arquivo de definicao seguidos de \enquote{entre outros}, na ordem do arquivo. A lista completa esta no repositorio, em \texttt{campos\_lexicais/catalogo\_termos\_aime.yaml}.}\label{tab:catalogo-etapa3}\\
\toprule
Campo & Procedencia & Amostra de termos do catalogo \\
\midrule
\endfirsthead
\caption[]{Catalogo lexical da Etapa 3 (AIME) (continuacao).}\\
\toprule
Campo & Procedencia & Amostra de termos do catalogo \\
\midrule
\endhead
\midrule
\multicolumn{3}{r@{}}{\footnotesize continua na proxima pagina}\\
\endfoot
\bottomrule
\endlastfoot
\texttt{modes\_of\_existence} & Quadro 1 de AIME / Conclusao de AIME & \texttt{mode of existence}, \texttt{modes of existence} \\
\texttt{preposition} & Quadro 1 de AIME / Conclusao de AIME & \texttt{preposition}, \texttt{prepositions}, \texttt{prepositional} \\
\texttt{felicity} & Quadro 1 de AIME / Conclusao de AIME & \texttt{felicity}, \texttt{infelicity}, \texttt{felicity condition}, \texttt{infelicity condition}, \texttt{felicity conditions}, \texttt{infelicity conditions}, \texttt{felicitous}, \texttt{infelicitous} \\
\texttt{diplomacy} & Quadro 1 de AIME / Conclusao de AIME & \texttt{diplomatic}, \texttt{diplomat}, \texttt{diplomats}, \texttt{diplomacy}, \texttt{negotiation}, \texttt{negotiations}, \texttt{negotiating}, \texttt{negotiate}, entre outros \\
\texttt{gaia} & Quadro 1 de AIME / Conclusao de AIME & \texttt{Gaia}, \texttt{Anthropocene}, \texttt{terrestrial}, \texttt{ecology}, \texttt{ecological}, \texttt{ecologies}, \texttt{climate}, \texttt{climatologist}, entre outros \\
\texttt{values} & Quadro 1 de AIME / Conclusao de AIME & \texttt{qualify values}, \texttt{qualifies values}, \texttt{qualifying values}, \texttt{qualified values}, \texttt{value system}, \texttt{values system}, \texttt{system of values}, \texttt{qualify} \\
\texttt{trajectory\_pass} & Quadro 1 de AIME / Conclusao de AIME & \texttt{trajectory}, \texttt{trajectories}, \texttt{course of action}, \texttt{hiatus}, \texttt{discontinuity}, \texttt{discontinuities}, \texttt{continuity}, \texttt{continuities}, entre outros \\
\texttt{institution} & Quadro 1 de AIME / Conclusao de AIME & \texttt{instituting}, \texttt{institute}, \texttt{institutes}, \texttt{instituted}, \texttt{institution}, \texttt{institutions} \\
\texttt{domain\_category} & Quadro 1 de AIME / Conclusao de AIME & \texttt{category mistake}, \texttt{category mistakes}, \texttt{domain}, \texttt{domains}, \texttt{crossing}, \texttt{crossings} \\
\texttt{moderns} & Quadro 1 de AIME / Conclusao de AIME & \texttt{Moderns}, \texttt{modernity}, \texttt{modernizing}, \texttt{modernism}, \texttt{modernization} \\
\texttt{alteration} & Quadro 1 de AIME / Conclusao de AIME & \texttt{alteration}, \texttt{alterations}, \texttt{alter}, \texttt{altered}, \texttt{alters}, \texttt{altering}, \texttt{being-as-other}, \texttt{being-as-another}, entre outros \\
\texttt{experience} & Quadro 1 de AIME / Conclusao de AIME & \texttt{experience}, \texttt{experiences}, \texttt{experienced}, \texttt{experiential}, \texttt{empirical} \\
\end{longtable}
